On considère l’équation différentielle (E) : $y' - y = 2(x + 1)\e^x$ où $y$ désigne une fonction de la variable réelle $x$ définie et dérivable sur $\R$, $y'$ sa fonction dérivée.

\begin{enumerate}
	\item Donner, puis résoudre l'équation homogène (E$_0$) :\sautp
	
	\Lignes{2}
	\item Vérifier que la fonction $p(x)=(x^2+2x)\e^{x}$ est une solution particulière de (E).\sautp

    \Lignes{8}
	%\item A l’aide du logiciel Géogébra, déterminer les nombres réels $a$ et $b$ pour que la fonction $p$ définie sur $\R$ par : $p(x) = (a~x^2+b~x)\e^x$ soit une solution particulière de l'équation différentielle (E), et donner l’expression de $p$. 

	
	%On trouve : \hfill $a=$\makebox[0.1\linewidth]{\dotfill} \hfill $b=$\makebox[0.1\linewidth]{\dotfill} \hfill $p(x)=$\makebox[0.3\linewidth]{\dotfill}

	%\smallskip
	
	%\appelprof
	
	\item En déduire l'ensemble des solutions de (E) : \sautp
	
	\Lignes{1}
	\item Déterminer, à l'aide de géogebra et d'un curseur bien choisi, la solution particulière $f$ de (E) telle que $f'(0)=3$.\sautp

	\Lignes{1}
    \appelprof
	%\item Déterminer, par un calcul, la solution $f$ de (E) telle que $f'(0)=3$ : 
	
	%{\it Commencer par calculer $f'(x)$, la dérivée, puis trouver $k$ et enfin répondre en donnant $f$.}\sautp
	
	%\Lignes{8} 

\end{enumerate}